\documentclass[a4paper,12pt]{article}
\pagestyle{plain}

\usepackage{indentfirst}
\usepackage{CJKutf8}
\usepackage[pdftex]{color,graphicx}

\begin{CJK}{UTF8}{gkai}
  \begin{document}
  \begin{center}
    \huge \bf  “送菜”需求及执行
  \end{center}


  \section{动机}
  而立过后，切盼有所成。 \\
  \begin{itemize}
  \item 一为争取时间

    如今上班终日无所事事，却不得不日日打卡报道； 家中大小事务，须先征得领导同意，方可抽身照顾。

  \item 二为获取金钱

    现工作虽养得人活，但长远看来，公司无用我之处，自然也谈不上发展前途。

  \item 三为实现自我

    十年费心所学，无以影响他人。 他日回首，只言浑噩度日，一事无成。

  \end{itemize}


  \section{原则}
  全心全意，做一件能影响他人的事，并从中获取回报。

  原则： \textbf{从善、勤奋、人性、睿智。}
  


  \section{解决的问题}
  帮用户解决生活中存在、且确实令其困扰的一个问题。

  综合自身兴趣、市场空间、成熟度、可行性几个方面的综合考虑。 决定选择“送菜”这个方向，
  解决用户如下问题：
  \begin{center}
    \includegraphics{need.png}
  \end{center}

  针对以上三点，提出如下几点要求，以作为工作中的准则，才能准确的满足用户需求：
  \begin{itemize}
  \item 专业的执行团队
  \item 严格的卫生标准
  \item 高校的服务质量
  \item 透明的投诉处理
  \end{itemize}

  \subsection{参考阅读}
  \begin{itemize}
  \item 行业分析

    派代讨论 http://bbs.paidai.com/topic.php?id=4365

    深圳行业分析 http://www.chinese-plant.com/hua/01-22/7063.htm

    望家欢访谈 http://hi.baidu.com/elnong/blog/item/4a2fc81113013d10b8127b7d.html

    高军：“我后来算了一笔账，其实家庭业务做的越多，亏损的越厉害。” 因为单个生意小，铺得摊子越大，赔得越厉害。 --- 进菜、等单、送菜(1小时内送达)的模式。

  \item 同行网站

    望家欢 http://www.wangjiahuan.net

    深圳家旺 http://www.szjjw.cn

    深圳禄苑 http://www.zgluyuan.com

    康尔信农产品 http://www.kangerxin.cn/

    田园送菜 http://www.tianyuansongcai.com/
  \end{itemize}

  
  \section{解决方案}
  通过网站、手机客户端\footnote{技术II期工程} ＋ 运营团队的方式为用户提供服务。 3分技术、7分运营。
  

  \section{大致需求}
  
  \subsection{基础功能}
  注册、登录、整站框架。
  
  \subsection{送}
  \begin{itemize}
  \item 选菜

    菜单分类： 晕、素传统分类，还是安营养成分（蛋白、维生素、脂肪等）分类？ \footnote{此处需要专业人士支持}
  \item 定菜

    提前多久？ 取消订单、确认收货、订单历史记录
  \item 收菜

    收菜地点、质量问题
  \item 充值
    
    支付宝、人工支付、积分支付、消费记录
  \end{itemize}

  \subsection{荐}
  \begin{itemize}
  \item 按地区当季菜推荐
  \item 用户分享推荐
  \item 病人、专题（特殊进货限量）食谱推荐 \footnote{此处需要专业人士支持}
  \end{itemize}

  \subsection{检}
  为本机构检验过的食品打上显眼绿标，可以放心购买。


  \section{执行步骤}
  \begin{itemize}
  \item 市场调查（需求分析、运作方法、资金规模、人力配置）
  \item 网站搭建
  \item 认同感评估
  \item 资金筹集
  \item 团队建设
  \item 深圳首发
  \end{itemize}

  
\end{CJK}
\end{document}
